\subsection{Independent checks of the spectral and algebraic formulas}

The following computations explain emendations in occurrences
1977-442.4, 1977-442.6--.7, 1977-443.9 and 1977-443.12--.14, and two later
algebraic formulas. They are checks made for this companion, with no claim
of historical novelty. The relevant published subjects are described by
\citet{calogerodegasperis1976inverse,calogerodegasperis1978extension,
calogero1978motion}; only their publisher metadata and abstracts or
editorial footnotes were available for the present comparison. Thus the
proofs below verify the displayed notebook equations independently of a
claim to have read their full published normalization conventions.

\begin{proposition}[Transmission phase in the notebook convention]\oursComputed
Let \(u\in C(\mathbb R)\cap L^1(\mathbb R)\) be real. Use the
notebook's physical scattering normalization for
\(\psi''=(u-k^2)\psi\): at negative infinity,
\(\psi\sim T(k)e^{-ikx}\), and at positive infinity,
\(\psi\sim e^{-ikx}+R(k)e^{ikx}\). With
\(T=|T|e^{2i\varphi}\) and the branch \(\varphi\to 0\) as
\(k\to+\infty\), its leading high-energy phase is
\[
 \varphi(k)=-\frac{1}{4k}\int_{\mathbb R}u(x)\,dx+o(k^{-1}).
\]
Consequently, the source's coefficient \(-c_0/2\), where
\(c_0=\int u\), is incompatible with its own factor 2 in the phase of T.
\end{proposition}
\begin{proof}
Normalize the left Jost solution f by \(f\sim e^{-ikx}\) at negative
infinity. Its Volterra equation is
\[
 f(x)=e^{-ikx}+\frac1k\int_{-\infty}^x
          \sin(k(x-y))u(y)f(y)\,dy.
\]
For real \(k>0\), Gronwall's inequality gives
\(\sup_x|f(x)|\le\exp(\|u\|_1/k)\) and
\(\sup_x|f(x)-e^{-ikx}|\le
\|u\|_1\exp(\|u\|_1/k)/k\).
Splitting the sine into exponentials shows that the coefficient of
\(e^{-ikx}\) at positive infinity is
\[
 A(k)=1-\frac1{2ik}\int_{\mathbb R}e^{iky}u(y)f(y)\,dy
     =1+\frac{i c_0}{2k}+O(k^{-2}).
\]
The error bound follows from the uniform estimate and integrability
of u. For all sufficiently large k, \(A(k)\ne0\).
The physical solution is \(\psi=f/A\), so \(T=A^{-1}\), and
\(\log T=-i c_0/(2k)+O(k^{-2})\). Taking half the imaginary part
on the branch near1 proves the formula. No convergence of a full
asymptotic series is asserted.

An independent exact check uses
\(u=-2r^2\operatorname{sech}^2(rx)\), \(r>0\), and
\[
 \psi=e^{-ikx}\frac{ik+r\tanh(rx)}{ik+r},\qquad
 T(k)=\frac{k+ir}{k-ir}.
\]
Direct differentiation gives \(\psi''=(u-k^2)\psi\). Its multiplying
factor tends to 1 at positive infinity and to
\((ik-r)/(ik+r)=(k+ir)/(k-ir)=T(k)\) at negative infinity,
so it has exactly the stated physical normalization, with \(R=0\).
Here \(c_0=-4r\) and \(\varphi=\arctan(r/k)=r/k+O(k^{-3})\),
again giving \(-c_0/4\). The symbolic residual and expansion are
recorded in the accompanying verification output. An independent
review caught an earlier reciprocal error in the stated boundary
normalization; these explicit endpoint checks correct it without
altering the final coefficient \(-c_0/4\).
\end{proof}

\begin{proposition}[Weights of the integrated recursion moments]\oursComputed
Let \(u\in\mathcal S(\mathbb R)\), and for
\(f\in\mathcal S(\mathbb R)\) define
\[
 L_u f=f''-4uf+2u'\int_x^\infty f(s)\,ds,\qquad
 S_u=xu'+2u,\qquad c_n[u]=\int_{\mathbb R}L_u^n S_u\,dx,
\]
The upper-tail primitive of a Schwartz function is bounded; its
product with \(u'\) is Schwartz. Thus L preserves this space,
all the displayed iterates and integrals exist, and the boundary
terms used below vanish.
For \(u_\lambda(x)=\lambda^2u(\lambda x)\), \(\lambda>0\),
\[
 c_n[u_\lambda]=\lambda^{2n+1}c_n[u].
\]
Thus a pure scaling flow \(u_t=b(xu_x+2u)\), for constant \(b\in\mathbb R\),
has \(\dot c_n=(2n+1)bc_n\), rather than \(nbc_n\).
\end{proposition}
\begin{proof}
If \(h_\lambda(x)=\lambda^a h(\lambda x)\), then its upper-tail
primitive has weight \(a-1\). Each of the three terms in
\(L_{u_\lambda}h_\lambda\) therefore has weight \(a+2\). Since
\(S_{u_\lambda}(x)=\lambda^2S_u(\lambda x)\), induction gives
weight \(2n+2\) before integration and \(2n+1\) afterwards.
The scaling-flow solution is obtained by \(\lambda=e^{bt}\).

Independently, integration by parts gives \(c_0=\int u\).
Put \(U(x)=\int_x^\infty u(s)\,ds\). Then
\(\int_x^\infty S_u=U-xu\), and the formula for L gives
\(c_1=-3\int u^2\). For example, on
\(u_\lambda=\lambda^2e^{-\lambda^2x^2}\), the two quantities are
\(\lambda\sqrt\pi\) and \(-3\lambda^3\sqrt{\pi/2}\).
These exact integrals independently test weights 1 and 3, already
refuting the source's proposed weights 0 and 1.
\end{proof}

\begin{proposition}[The sign in the Bianchi quadrilateral]\oursComputed
Let \(p,q\in\mathbb C\) be constants, and let U,A,B be
complex-valued \(C^1\) functions on a real interval I. Suppose the
two lower B\"acklund edges satisfy
\[
 U_x+A_x=-\tfrac12(A-U)(4p+A-U),
\]
and the same equation with \((A,p)\) replaced by \((B,q)\).
On any subinterval where the displayed denominator is nonzero,
the compatible upper value is the \(C^1\) function
\[
 C=U-\frac{(p+q)(A-B)}{p-q+\tfrac12(A-B)}.
\]
In particular, the denominator uses \(A-B\), with the same order
as the numerator in these conventions.
\end{proposition}
\begin{proof}
The upper edges require
\[
 A_x+C_x=-\tfrac12(C-A)(4q+C-A),\qquad
 B_x+C_x=-\tfrac12(C-B)(4p+C-B).
\]
Subtracting them and eliminating \(A_x-B_x\) by the two lower
edges gives a linear equation for C. Solving it yields the displayed
formula. This elimination is one check.

For a second check, differentiate that rational formula using the two
lower equations, leaving \(U_x\) arbitrary. Substitution into either
upper equation gives zero identically. The symbolic verification file
performs these as two separate residual calculations. The computation
only concerns this local differential quadrilateral; it neither fills
in missing scattering norming data nor proves a global permutability
statement at zeros of the denominator.
\end{proof}

\begin{proposition}[The Hermite interaction matrix]\oursComputed
Let \(\bar x_1,\ldots,\bar x_n\) be the distinct zeros of the
physicists' Hermite polynomial \(H_n\), and put
\[
 A_{jk}=-\frac1{(\bar x_j-\bar x_k)^2}\quad(j\ne k),\qquad
 A_{jj}=\sum_{k\ne j}\frac1{(\bar x_j-\bar x_k)^2}.
\]
For the source's root dynamics
\[
 \ddot x_j=-x_j+\sum_{k\ne j}
       \frac{1+2\dot x_j\dot x_k}{x_j-x_k},
\]
the equilibrium linearization is
\(\ddot\varepsilon+(I+A)\varepsilon=0\).
Moreover, \(\operatorname{Spec}(A)=\{0,1,\ldots,n-1\}\), with
each eigenvalue simple.
\end{proposition}
\begin{proof}
The derivative of each inverse difference contributes
\(-(\varepsilon_j-\varepsilon_k)/(\bar x_j-\bar x_k)^2\).
The velocity products have zero linear part at a stationary point.
The independent term \(-x_j\) contributes \(-\varepsilon_j\),
giving the stated matrix \(I+A\).

For the spectrum, use the notebook polynomial equation
\[
 \Psi_{tt}+\tfrac12\Psi_{xx}-x\Psi_x+n\Psi=0,
 \qquad
 \Psi=2^{-n}\sum_{m=0}^n c_m(t)H_{n-m}(x),\quad c_0=1.
\]
The Hermite differential equation gives
\(\ddot c_m+mc_m=0\) for \(1\le m\le n\).
At a polynomial with distinct roots, the differential of the
root-to-coefficient map is invertible: a variation has
\(\delta\Psi(\bar x_j)=-\Psi'(\bar x_j)\delta x_j\), and
evaluation at n distinct points is invertible on polynomials of
degree at most \(n-1\). Therefore \(I+A\) is conjugate to
\(\operatorname{diag}(1,\ldots,n)\).

As a second explicit check, for n=3 the roots are
\(-\sqrt{3/2},0,\sqrt{3/2}\) and
\[
 A=\begin{pmatrix}
 5/6&-2/3&-1/6\\
 -2/3&4/3&-2/3\\
 -1/6&-2/3&5/6
 \end{pmatrix},\qquad
 \det(\lambda I-A)=\lambda(\lambda-1)(\lambda-2).
\]
The verification code independently differentiates the full nonlinear
force and obtains its Jacobian \(-I-A\). It also checks the Hermite
mode identity exactly for all 36 pairs \(1\le m\le n\le8\).
Already n=1 shows why the identity cannot be omitted: A is zero but
the single root satisfies \(\ddot x=-x\).
\end{proof}

\begin{remark}[Two later algebraic checks]\oursComputed
For occurrence1977-457.18, take the ordinary \(\mathfrak{sl}_2\)
relations and define
\[
 \varepsilon_{3r}=\tfrac12ht^r,\quad
 \varepsilon_{3r+1}=et^r,\quad
 \varepsilon_{3r-1}=\tfrac12ft^r.
\]
Matrix commutators verify all nine residue cases of
\([\varepsilon_i,\varepsilon_j]=a_{ij}\varepsilon_{i+j}\),
where \(a_{ij}\in\{-1,0,1\}\) is congruent to \(j-i\) modulo 3.
Separately, the 27 residue triples satisfy Jacobi directly. The
nonhomogeneous ideal \((t-1)\mathfrak{sl}_2[t,t^{-1}]\) explains why
the classification convention is graded simplicity, as already stated
in \citet[Introduction]{kac1968graded}.

For occurrence1977-457.25, left Grassmann differentiation satisfies
\(\{\partial_{\theta^a},\theta^b\}=\delta_a^b\), while distinct
odd multiplications or derivatives anticommute. Hence the two cross
terms in the notebook's Q and \(\bar Q\) each contribute
\(-i\sigma^m\partial_m\), giving
\(\{Q_\alpha,\bar Q_{\dot\beta}\}=-2i\sigma^m_{\alpha\dot\beta}
\partial_m\); the same-chirality brackets vanish. Independently,
the code realizes multiplication and differentiation as exact 16-by-16
matrices on the exterior algebra of four generators and checks all 12
anticommutator pairs with four independent commuting derivative
symbols. These calculations validate the surviving manuscript signs;
they do not resolve the illegible cancelled mark or supply the missing
Lorentz action on the spinor coordinates.
\end{remark}
